documentclass[a4paper,12pt]{article}

usepackage{graphicx}
usepackage{geometry}
usepackage{tikz}
usetikzlibrary{shapes,arrows.meta,positioning,calc}
usepackage{booktabs}
usepackage{caption}
usepackage{float}
usepackage{hyperref}

geometry{a4paper, top=2.7cm, bottom=2.7cm, left=2.5cm, right=2.5cm}
graphicspath{{images}}

usepackage{xepersian}
settextfont{XB Niloofar}

title{textbf{گزارش آزمایش اول آشنایی با محیط‌های شبیه‌سازی مدارهای منطقی} [4pt]
large آشنایی با نرم‌افزارهای Fritzing، Logisim و Proteus}
author{علی اصغر طایری - ۴۰۴۱۷۱۱۳۳  رضا حضرتی - ۴۰۴۱۰۵۷۶۷}
date{درس آزمایشگاه مدارهای منطقی}

begin{document}

maketitle
thispagestyle{empty}
newpage

tableofcontents
newpage

section{مقدمه و تعریف مسئله}
هدف از این آزمایش، آشنایی عملی با سه محیط نرم‌افزاری پرکاربرد در طراحی و شبیه‌سازی مدارهای منطقی و الکترونیکی است textbf{Fritzing}، textbf{Logisim} و textbf{Proteus}. این سه ابزار در سه سطح متفاوت از انتزاع کار می‌کنند Fritzing برای بستن مدارهای ساده روی بردبورد و مشاهدهٔ اتصالات فیزیکی، Logisim برای طراحی و شبیه‌سازی مدار در سطح گیت‌های منطقی، و Proteus برای طراحی مدارهای ترکیبی پیچیده‌تر با نمای شماتیک حرفه‌ای و امکان شبیه‌سازی نزدیک به واقعیت.

مسئله‌ای که در این آزمایش باید حل شود در سه بخش خلاصه می‌شود
begin{enumerate}
item بستن یک مدار ساده و سپس یک مدار مبتنی بر تراشهٔ lr{7404} روی بردبورد در lr{Fritzing} و ترسیم اتصالات داخلی آن.
item ساخت مدار جمع‌کنندهٔ کامل (lr{Full Adder}) در lr{Logisim} بر اساس جدول درستی آن، و سپس گسترش آن به یک مدار جمع‌کنندهتفریق‌کنندهٔ چهار بیتی با استفاده از سیگنال کنترلی $C_{in}$.
item پیاده‌سازی یک جمع‌کنندهٔ چهاربیتی از نوع lr{Carry-Look-Ahead} در نرم‌افزار lr{Proteus}.
end{enumerate}

هدف نهایی، درک تفاوت این سه سطح شبیه‌سازی و کسب توانایی لازم برای استفاده از آن‌ها در آزمایش‌های آینده است.

section{بررسی و انتخاب روش طراحی}

subsection{بخش Fritzing}
از آنجا که هدف این بخش آشنایی با نحوهٔ اتصال فیزیکی قطعات روی بردبورد است، مدارها ابتدا در نمای lr{Breadboard} بسته شدند و سپس با نمای lr{Schematic} مقایسه گردیدند تا ارتباط بین دو نما درک شود. مدار اول (شکل ۱) شامل یک مقاومت $220,Omega$، یک lr{LED} قرمز و یک باتری است. مدار دوم با استفاده از تراشهٔ lr{7404} (شش گیت lr{NOT}) بسته شد؛ با استفاده از 6 گیت lr{NOT} به‌صورت زنجیره‌ای و بازخورد آن به ورودی، یک مدار نوسان‌ساز حلقه‌ای (lr{Ring Oscillator}) ساخته شد که چشمک‌زدن lr{LED} نتیجهٔ آن است.

subsection{بخش Logisim}
برای ساخت جمع‌کنندهٔ کامل، به‌جای طراحی آزمایشی، از فرم استاندارد جبر بول بر اساس جدول درستی استفاده شد
[
S = a oplus b oplus C_{in}, qquad C_{out} = ab + C_{in}(a oplus b)
]
این انتخاب به این دلیل صورت گرفت که فرم فوق کمترین تعداد گیت (۲ گیت lr{XOR}، ۲ گیت lr{AND} و ۱ گیت lr{OR}) را برای پیاده‌سازی جمع‌کنندهٔ کامل نیاز دارد.

برای گسترش به مدار جمعتفریق چهاربیتی، چهار جمع‌کنندهٔ کامل به‌صورت آبشاری (lr{Ripple Carry}) به هم متصل شدند. برای پشتیبانی از عملیات تفریق، از روش lr{2's Complement} استفاده شد ورودی $B$ پیش از ورود به هر جمع‌کننده از یک گیت lr{XOR} با سیگنال کنترلی $C_{in}$ عبور داده می‌شود (که هم‌زمان به‌عنوان کری ورودی مرحلهٔ اول نیز استفاده می‌شود). وقتی $C_{in}=0$ عملیات $F=A+B$ و وقتی $C_{in}=1$ عملیات $F=A-B$ انجام می‌شود.

subsection{بخش Proteus}
برای این بخش، طراحی مستقیماً بر اساس معادلات lr{Carry-Look-Ahead} در سطح گیت انجام شد و از تراشهٔ آماده استفاده نشد. برای هر بیت، سیگنال‌های $P_i=A_ioplus B_i$ با گیت‌های lr{XOR} و $G_i=A_iB_i$ با گیت‌های lr{AND} تولید شدند؛ سپس با بسط رابطهٔ $C_{i+1}=G_i+P_iC_i$ برای بیت‌های بالاتر، هر کری به‌صورت مستقیم و موازی از ترکیب گیت‌های lr{AND} و lr{OR} روی $P$ و $G$ مراحل قبل محاسبه شد (بدون انتظار برای انتشار کری از بیت‌های پایین‌تر). در نهایت خروجی مجموع هر بیت با یک گیت lr{XOR} از $P_i$ و کری ورودی همان بیت به دست آمد. این انتخاب به قیمت افزایش تعداد گیت‌ها (مطابق شکل مدار، مجموعاً ۲۸ گیت)، سرعت محاسبهٔ کری را نسبت به روش آبشاری بهبود می‌دهد.


section{بلوک‌دیاگرام طرح پیشنهادی اولیه}

پیش از ورود به جزئیات مدار، بلوک‌دیاگرام کلی هر بخش به شکل زیر در نظر گرفته شد.
begin{figure}[H]
centering
begin{tikzpicture}[
  box.style={draw, thick, minimum width=2.6cm, minimum height=2cm, align=center},
  port.style={font=small},
  bus.style={-, thick},
  =Latex,
  node distance=1.6cm
]

node[box] (fa) {Full Adder};

node[port, left=1.8cm of fa.west, yshift=0.5cm]  (a)   {$a$};
node[port, left=1.8cm of fa.west, yshift=0cm]    (b)   {$b$};
node[port, left=1.8cm of fa.west, yshift=-0.5cm] (cin) {$C_{in}$};

node[port, right=1.6cm of fa.east, yshift=0.5cm]  (s)    {$S$};
node[port, right=1.6cm of fa.east, yshift=-0.5cm] (cout) {$C_{out}$};

draw[bus] (a.east)   -- (fa.west - a.east);
draw[bus] (b.east)   -- (fa.west - b.east);
draw[bus] (cin.east) -- (fa.west - cin.east);
draw[bus] (fa.east - s.west)    -- (s.west);
draw[bus] (fa.east - cout.west) -- (cout.west);

begin{scope}[xshift=8cm]
node[box, minimum width=3cm] (as) {AdderSubtractor};

node[port, left=1.8cm of as.west, yshift=0.7cm]  (A)   {$A_{30}$};
node[port, left=1.8cm of as.west, yshift=0cm]    (B)   {$B_{30}$};
node[port, left=1.8cm of as.west, yshift=-0.7cm] (Cin) {$C_{in}$};

node[port, right=1.6cm of as.east, yshift=0.45cm]  (F)    {$F_{30}$};
node[port, right=1.6cm of as.east, yshift=-0.45cm] (Cout) {$C_{out}$};

draw[bus] (A.east)   -- node[midway, above, font=scriptsize] {4} (as.west - A.east);
draw[bus] (B.east)   -- node[midway, above, font=scriptsize] {4} (as.west - B.east);
draw[bus] (Cin.east) -- (as.west - Cin.east);
draw[bus] (as.east - F.west)    -- node[midway, above, font=scriptsize] {4} (F.west);
draw[bus] (as.east - Cout.west) -- (Cout.west);
end{scope}

end{tikzpicture}
caption{بلوک‌دیاگرام پیشنهادی جمع‌کنندهٔ کامل (چپ) و مدار جمعتفریق چهاربیتی (راست)}
end{figure}
مدار Fritzing نیز در سطح بلوک به‌صورت زیر در نظر گرفته شد
begin{itemize}
    item منبع تغذیه $rightarrow$ مقاومت محدودکنندهٔ جریان $rightarrow$ lr{LED} (برای مدار اول)
    item منبع تغذیه $rightarrow$ زنجیرهٔ گیت‌های lr{NOT} با بازخورد $rightarrow$ lr{LED} خروجی (برای مدار دوم)
end{itemize}
section{روند پیاده‌سازی}

subsection{پیاده‌سازی در Fritzing}
ابتدا یک بردبورد خالی در نظر گرفته شد و با استفاده از ماوس، نحوهٔ اتصال داخلی ردیف‌ها و ریل‌های تغذیهزمین بررسی و رنگ‌آمیزی شد تا ساختار درونی بردبورد شناخته شود (شکل ۲). سپس در نمای بردبورد، مقاومت $220,Omega$ و lr{LED} قرمز با پایانه‌های باتری مطابق شکل استاندارد به هم متصل شدند (شکل ۳). سپس، تراشهٔ lr{7404} روی بردبورد قرار گرفت (شکل ۴) و با اتصال زنجیره‌ای سه گیت lr{NOT} و بازخورد خروجی آخرین گیت به ورودی گیت اول، مدار نوسان‌ساز حلقه‌ای بسته شد.

begin{figure}[H]
centering
begin{minipage}{0.45linewidth}
    centering
    includegraphics[width=linewidth]{1-2-1-1.png}
end{minipage}hfill
begin{minipage}{0.45linewidth}
    centering
    includegraphics[width=linewidth]{1-2-1-2.png}
end{minipage}

vspace{0.4cm}

begin{minipage}{0.45linewidth}
    centering
    includegraphics[width=linewidth]{1-2-1-3.png}
end{minipage}hfill
begin{minipage}{0.45linewidth}
    centering
    includegraphics[width=linewidth]{1-2-1-4.png}
end{minipage}
caption{بررسی اتصالات داخلی بردبورد}
end{figure}

begin{figure}[H]
centering
includegraphics[width=0.45linewidth]{1-2-2-1.png}hspace{0.3cm}
includegraphics[width=0.45linewidth]{1-2-2-2.png}
caption{مدار ساده مقاومت-LED-باتری بسته‌شده روی بردبورد}
end{figure}

begin{figure}[H]
centering
includegraphics[width=0.62linewidth]{{Screenshot 2026-07-30 063940}.png}
caption{قرارگیری تراشهٔ lr{7404} روی بردبورد پیش از سیم‌کشی نهایی}
end{figure}

begin{figure}[H]
centering
includegraphics[width=0.75linewidth]{{Screenshot 2026-08-04 160840}.png}
caption{نمای نهایی مدار نوسان‌ساز حلقه‌ای در نرم‌افزار Fritzing}
end{figure}

subsection{پیاده‌سازی در Logisim}
در مرحلهٔ اول، جمع‌کنندهٔ کامل با دو گیت lr{XOR} (برای تولید $S$)، دو گیت lr{AND} و یک گیت lr{OR} (برای تولید $C_{out}$) ساخته شد (شکل 6). سپس هر ۸ حالت ممکن ورودی $(a,b,C_{in})$ اعمال و خروجی‌ها با جدول درستی مقایسه شد.

در مرحلهٔ دوم، چهار نمونه از این جمع‌کنندهٔ کامل به‌صورت آبشاری به هم متصل شدند و در مسیر ورودی $B$ هر مرحله یک گیت lr{XOR} با سیگنال کنترلی مشترک $C_{in}$ قرار گرفت. کری خروجی هر مرحله به کری ورودی مرحلهٔ بعد وصل شد و کری ورودی مرحلهٔ اول نیز به همان سیگنال کنترلی $C_{in}$ متصل گردید تا مدار در حالت $C_{in}=1$ عمل تفریق (متمم دوی $B$ به‌علاوهٔ یک) را انجام دهد.

begin{figure}[H]
centering
includegraphics[width=0.6linewidth]{1.PNG}
caption{مدار جمع‌کنندهٔ کامل ساخته‌شده در Logisim (حالت $a=b=C_{in}=0$)}
end{figure}

begin{figure}[H]
centering
includegraphics[width=0.55linewidth]{4444.PNG}
caption{مدار نهایی جمعتفریق چهاربیتی در Logisim -- حالت جمع ($C_{in}=0$)}
end{figure}

begin{figure}[H]
centering
includegraphics[width=0.55linewidth]{44444.PNG}
caption{مدار نهایی جمعتفریق چهاربیتی در Logisim -- حالت تفریق ($C_{in}=1$)}
end{figure}

subsection{پیاده‌سازی در Proteus}
برخلاف روش آبشاری، در این پیاده‌سازی به‌جای استفاده از تراشه‌های آماده، مدار CLA کاملاً از گیت‌های منطقی مجزا ساخته شد ابتدا با گیت‌های lr{XOR} سیگنال‌های lr{Propagate} ($P_i=A_ioplus B_i$) و با گیت‌های lr{AND} سیگنال‌های lr{Generate} ($G_i=A_iB_i$) برای هر بیت به‌طور مستقل تولید شدند. سپس با ترکیب چند مرحله‌ای گیت‌های lr{AND} و lr{OR} بر اساس رابطهٔ $C_{i+1}=G_i+P_iC_i$، کری هر بیت مستقیماً از روی $P$ و $G$ مراحل قبل محاسبه گردید. در نهایت خروجی مجموع هر بیت با یک گیت lr{XOR} از ترکیب $P_i$ و کری ورودی همان بیت به دست آمد.

begin{figure}[H]
centering
includegraphics[width=0.75linewidth]{{Screenshot 2026-07-30 080411}.png}
caption{مدار نهایی جمع‌کنندهٔ چهاربیتی lr{Carry-Look-Ahead} در Proteus}
end{figure}


section{دیاگرام یا شماتیک طرح اولیه}

شماتیک اولیه‌ای که در ابتدای آزمایش برای پیاده‌سازی پیشنهاد شده بود، شامل مدار مقاومت-LED (شکل ۱ صورت آزمایش)، مدار نوسان‌ساز مبتنی بر lr{7404} (شکل ۲ صورت آزمایش)، جمع‌کنندهٔ کامل با جدول درستی استاندارد (شکل ۳ صورت آزمایش)، بلوک جمعتفریق چهاربیتی (شکل ۴ صورت آزمایش) و ساختار جمع‌کنندهٔ lr{Carry-Look-Ahead} (شکل ۵ صورت آزمایش) بود. این طرح اولیه بدون تغییر ساختاری اجرا شد و تنها در انتخاب مقادیر جزئی (مانند رنگ lr{LED} یا آرایش فیزیکی قطعات روی بردبورد) با توجه به قطعات در دسترس، تغییرات جزئی اعمال گردید.

section{بررسی Case Test ها}

subsection{تست جمع‌کنندهٔ کامل}
هر ۸ حالت ممکن ورودی به مدار جمع‌کنندهٔ کامل در Logisim اعمال شد و خروجی با جدول درستی مرجع مقایسه گردید. نتیجه در جدول زیر خلاصه شده است.
begin{table}[H]
centering
caption{جدول درستی جمع‌کنندهٔ کامل و نتیجهٔ تست}
begin{LTR}
begin{tabular}{ccccccc}
toprule
multicolumn{4}{c}{ } & multicolumn{2}{c}{خروجی} & 
$a$ & $b$ & $C_{in}$ & & $S$ & $C_{out}$ & نتیجه 
midrule
0 & 0 & 0 & & 0 & 0 & lr{OK} 
0 & 0 & 1 & & 1 & 0 & lr{OK} 
0 & 1 & 0 & & 1 & 0 & lr{OK} 
0 & 1 & 1 & & 0 & 1 & lr{OK} 
1 & 0 & 0 & & 1 & 0 & lr{OK} 
1 & 0 & 1 & & 0 & 1 & lr{OK} 
1 & 1 & 0 & & 0 & 1 & lr{OK} 
1 & 1 & 1 & & 1 & 1 & lr{OK} 
bottomrule
end{tabular}
end{LTR}
end{table}
در تمام ۸ حالت، خروجی مدار ساخته‌شده در Logisim دقیقاً با جدول درستی مرجع مطابقت داشت.

begin{figure}[H]
centering
includegraphics[width=0.4linewidth]{3.PNG}hspace{0.3cm}
includegraphics[width=0.4linewidth]{6.PNG}hspace{0.3cm}
includegraphics[width=0.4linewidth]{2.PNG}
caption{نمونه‌ای از تست سه حالت ورودی روی جمع‌کنندهٔ کامل}
end{figure}

subsection{تست مدار جمعتفریق چهاربیتی}
برای بررسی صحت عملکرد، به‌عنوان نمونه مقادیر $A=1101_2;(13)$ و $B=1011_2;(11)$ در نظر گرفته شد
begin{itemize}
item با $C_{in}=0$ خروجی $F = 1000_2 = 8$و یک کری ، که برابر $A+B$ است.
item با $C_{in}=1$ خروجی $F = 0010_2 = 2$، که برابر $A-B$ است.
end{itemize}
هر دو حالت با محاسبهٔ دستی مطابقت داشت که صحت پیاده‌سازی مدار جمعتفریق را تأیید می‌کند.

subsection{تست جمع‌کنندهٔ CLA}
مقادیر مختلف ورودی به مدار CLA در Proteus اعمال شد و مشاهده شد که کری هر بیت بلافاصله و بدون تأخیر انتشار از بیت‌های پایین‌تر محاسبه می‌شود؛ نتیجهٔ مجموع در تمامی حالت‌های آزموده‌شده با محاسبهٔ دستی یکسان بود.

begin{figure}[H]
centering
includegraphics[width=0.75linewidth]{{Screenshot 2026-08-02 073010}.png}
caption{نمونه‌ای از تست مدار جمع‌کنندهٔ lr{CLA} در Proteus}
end{figure}

begin{figure}[H]
centering
includegraphics[width=0.75linewidth]{{Screenshot 2026-08-02 072952}.png}
caption{نمونه‌ای دیگر از تست مدار جمع‌کنندهٔ lr{CLA} در Proteus}
end{figure}

section{نتیجه‌گیری}

به‌طور کلی، این آزمایش نشان داد که انتخاب ابزار شبیه‌سازی باید متناسب با سطح انتزاع مورد نیاز مسئله (فیزیکی، منطقی یا الکتریکی) صورت گیرد، و درک صحیح جمع‌کنندهٔ کامل به‌عنوان بلوک پایه، امکان ساخت مدارهای حسابی پیچیده‌تر مانند جمعتفریق‌کننده و CLA را فراهم می‌سازد.

end{document}